\section{Discussion}
\label{sec:discussion}

\myparagraph{Practical use} What the model does well is order tests, not decide whether each one is flaky. When it flags a test as flaky, it is right only 2 or 3 times out of 100, so using it to block a delivery in the pipeline would raise too many false alarms. Its usefulness appears when we consider a team that lacks the time to rerun every test and must choose where to start. In that case, the ranking works as a priority list: the top quarter of each project's list concentrates 77\% of the flaky tests (Section~\ref{sec:rq4}), and rerunning that quarter would cost about a quarter of the total effort, if each test took roughly the same time. An example helps show the gain: in a repository of 1,000 tests at the natural rate of 0.39\%, there are about 4 flaky tests; the top quartile has 250 tests and, with 77\% coverage, about 3 of them, a density of roughly 1.2\%, three times that of the whole set, though still low for gating. \Den{Um lift de 3$\times$ é modesto e precisa ser comparado com a heurística trivial ``rerodar só os testes com \texttt{task}/\texttt{async}'' (ver o comentário na RQ4; requer execução). Se a heurística der o mesmo ganho, a recomendação prática deveria ser ela, e o modelo fica dispensável. Isso precisa ser enfrentado explicitamente.} This holds as long as rerunning actually reveals the flakiness, which depends on how often each test fails. There is also a secondary benefit: the vocabulary analysis shows which constructs tend to accompany flakiness, such as fixed waits (\texttt{Sleep}), cancellation, and garbage collection, and a code reviewer can use these cues as warnings during review. We did not test this use with practitioners. The ranking can also select which tests to hand to dynamic tools that reproduce concurrency failures, such as FlakeRepro \cite{leesatapornwongsa2022flakerepro}, which need the test and an observed failure to work.

\myparagraph{Part of the signal is project-specific} When the model is trained and tested on tests from all projects mixed together, it performs very well (AUC of 0.86--0.87). When it must predict the tests of a project it has never seen, performance drops (AUC of 0.75--0.78). This drop indicates that part of what the model learns is each project's writing habits, not flakiness itself. Three repositories hold two thirds of the flaky tests (20 of 30), removing attributes reduces MCC by about 0.05 to 0.07, and the tokens with the highest information gain come largely from a single project. On the other hand, cross-repository performance remains well above chance, which shows that a general signal also exists. In practice, this calls for caution when importing results from other teams and languages, and, if the tool is adopted in an organization, it is worth calibrating it with the flaky tests that organization already knows about. \Den{Esta é a discussão central do paper de Go (``what classifiers actually learn''). Para manter a coerência entre os trabalhos do grupo, e porque o revisor pode ser o mesmo, adotar aqui os mesmos diagnósticos: baseline só com o nome do repositório e IG leave-one-project-out (requer execução). Sem eles, ``part of the signal is project-specific'' fica só qualitativo.}

\myparagraph{Comparison with other ecosystems} In Swift, the same kind of classifier reached an MCC of 0.75 under balanced cross-validation, against 0.08 for the best baseline \cite{medeiros2026vocabulary}; in C\#, the MCC is 0.51--0.53, and the non-trivial baselines reach 0.36--0.40. The precision, recall, and F1 of our best model are also lower than those reported for Java, Python, JavaScript, and Swift (Table~\ref{tab:literatura}). The numbers are not directly comparable, because the datasets, labels, protocols, and number of positives differ (91 flaky in Swift, 30 here), but they suggest that the flakiness vocabulary in C\# is less distinctive: asynchronous constructs are also common in stable tests, and there is no stability marker like Swift's \texttt{xctassertequal}.

\myparagraph{Toward better features} The errors point to natural extensions: propagating tokens through the helper methods called (interprocedural analysis) and designing features that separate fixed waits from controlled time. A simple version of the first, which adds to each test the tokens of the helpers it calls, brought small and inconsistent gains (Section~\ref{sec:sensitivity}); a proper analysis, with type information, is left for future work. The second requires distinguishing, in the code, a \texttt{Sleep} that assumes a duration from a \texttt{Delay} controlled by a simulated clock, which vocabulary alone cannot do.

\myparagraph{Methodological lessons} First, the unit of analysis changes the answer. A single test method can generate several cases that the runner reports separately: counting cases, there are 50 flaky tests; counting the methods that declare them, there are 30. In StreamJsonRpc, for example, the same test runs with six serializers and generates 25 flaky cases from only 8 methods. Second, labels depend on the round of runs: of the 36 flaky units found by either of the two rounds, only 13 are flaky in both, and the comparison between vocabulary and structure changes with the round (Section~\ref{sec:rq12}), so studies with a single round should report this sensitivity. Third, we document .NET-specific pitfalls that distorted our initial counts (Table~\ref{tab:armadilhas}): parameterized test names containing timestamps, thread ids, or dates, which make the same test look like several; differences between target frameworks; tests declared in base classes and executed in several derived classes; and contaminated environments, which produced spurious flaky tests.

\myparagraph{Recommendations for rerun studies in .NET} The pitfalls we encountered suggest a minimal protocol for anyone wishing to repeat this kind of study: (i) identify each test by its fully qualified class plus its display name and normalize timestamps, thread ids, and dates in the names; (ii) resolve inherited and parameterized tests to the declaring method and report both counts, per case and per method; (iii) build once, then run without rebuilding and on the same target framework in every round; (iv) isolate the environment between runs and discard contaminated rounds; and (v) report the sensitivity of the labels to the number of runs. \Den{Estas lições e este protocolo talvez sejam a contribuição mais original e reutilizável do artigo. Considerar uma seção própria (ou uma RQ), com o script de normalização/resolução de herança no artefato, como ferramenta. Seria ótimo quantificar quanto cada armadilha distorceria as contagens se não fosse tratada (ex.: quantos flaky espúrios surgem sem a normalização de nomes).}
